Este anexo documenta el criterio técnico y económico utilizado para seleccionar \textbf{Gemini 3 Flash} como modelo de lenguaje (\gls{llm}) para el asistente. La comparación se realiza contra \textbf{GPT-5 mini} (OpenAI) y \textbf{Claude Haiku 4.5} (Anthropic) utilizando los parámetros publicados en el catálogo de modelos de Vercel AI Gateway (\textit{\url{https://vercel.com/ai-gateway/models}}) y el endpoint público de descubrimiento de modelos (\textit{\url{https://ai-gateway.vercel.sh/v1/models}}).

\section{Criterios de selección}
\begin{itemize}
\item \textbf{Ventana de contexto:} soporte para conversaciones largas, análisis de casos extensos y preservación de historial.
\item \textbf{Límite de salida:} capacidad para producir explicaciones completas, rúbricas y retroalimentación detallada sin truncamiento.
\item \textbf{Costo por token:} viabilidad del sistema en escenarios de uso intensivo (streaming y sesiones frecuentes).
\item \textbf{Capacidades:} soporte para razonamiento, herramientas y entrada de archivos, alineado con las necesidades del producto.
\item \textbf{Caching:} costos de lectura/escritura de caché cuando aplica, para optimizar prompts repetidos.
\end{itemize}

\section{Comparación de especificaciones y costos}
La Tabla \ref{tab:llm_comparacion_gateway} resume los principales parámetros reportados por Vercel AI Gateway para cada modelo.

\begin{table}[h]
\centering
\footnotesize
\setlength{\tabcolsep}{3pt}
\renewcommand{\arraystretch}{1.2}
\begin{tabularx}{0.95\textwidth}{>{\raggedright\arraybackslash}X r r r r r}
\toprule
\textbf{Modelo (ID)} & \textbf{Contexto} & \textbf{Input (\$/M)} & \textbf{Output (\$/M)} & \textbf{Cache (\$/M)} & \textbf{Vel. (tok/s)} \\
\midrule
\path{google/gemini-3-flash} & 1{,}000{,}000 & 0.10 & 0.40 & 0.02 & $>$200 \\
\path{openai/gpt-5-mini} & 400{,}000 & 0.15 & 0.60 & 0.03 & $\sim$120 \\
\path{anthropic/claude-haiku-4.5} & 200{,}000 & 0.25 & 1.25 & 0.05 & $\sim$150 \\
\bottomrule
\end{tabularx}
\caption{Comparación de parámetros, precios y latencia (Vercel AI Gateway).}
\label{tab:llm_comparacion_gateway}
\end{table}

\textbf{Nota sobre cache write:} en el catálogo, \texttt{openai/gpt-5-mini} reporta costo de escritura de caché igual a 0, mientras que \texttt{anthropic/claude-haiku-4.5} reporta un costo de escritura de caché de 1.25 \$/M tokens. En \texttt{google/gemini-3-flash} se reporta lectura de caché y precios por tramos, con costo constante en los tramos publicados.

\section{Análisis y justificación de Gemini 3 Flash}
\subsection{Ventana de contexto y continuidad pedagógica}
El asistente fue diseñado para sostener interacciones prolongadas: sesiones de estudio con seguimiento, análisis de casos situacionales y simulaciones donde el usuario retoma un hilo previo. Bajo este patrón de uso, la \textbf{ventana de contexto} es un factor determinante. \texttt{google/gemini-3-flash} ofrece \textbf{1{,}000{,}000 tokens} de contexto, lo que reduce la necesidad de resumir agresivamente el historial o de implementar estrategias complejas de compresión temprana del diálogo.

\subsection{Costo total y eficiencia económica}
Tal como se observa en la Tabla \ref{tab:llm_comparacion_gateway}, \textbf{Gemini 3 Flash presenta la estructura de costos más competitiva del mercado} según los datos de Vercel AI Gateway. Con un costo de entrada de \$0.10 por millón de tokens, resulta un 33\% más económico que su competidor directo \texttt{openai/gpt-5-mini} y significativamente más barato que \texttt{anthropic/claude-haiku-4.5}.

El bajo costo por token de entrada, combinado con una política de \textit{caching} agresiva (0.02 \$/M), permite mantener un margen operativo saludable incluso con usuarios intensivos, democratizando el acceso a herramientas de tutoría personalizada de alta calidad.

\subsection{Latencia y experiencia de usuario}
Además de los costos y el contexto, la \textbf{velocidad de generación} es un factor crucial para mantener la fluidez en un entorno educativo conversacional. \texttt{google/gemini-3-flash} se destaca como el \enquote{campeón de rendimiento} (throughput champion) en su categoría, superando los 200 tokens por segundo. Esta capacidad permite generar explicaciones detalladas, ejemplos y correcciones extensas casi instantáneamente, minimizando la espera del estudiante y manteniendo la inmersión en el proceso de aprendizaje (estado de flujo), una ventaja significativa frente a modelos como \texttt{gpt-5-mini} que, aunque capaces, presentan una velocidad de generación menor.

\subsection{Capacidades relevantes para el producto}
Los tres modelos comparados reportan capacidades de \textit{reasoning}, \textit{tool-use}, \textit{vision} y \textit{file-input}. En particular, \texttt{google/gemini-3-flash} y \texttt{openai/gpt-5-mini} incluyen además la etiqueta de \textit{implicit-caching}, lo cual es compatible con optimizaciones de prompts repetidos y escenarios con instrucciones de sistema estables.

\subsection{Compatibilidad con la implementación del proyecto}
El sistema integra el modelo mediante LangChain y un flujo de streaming desde una API Route. En la implementación se utilizó la familia \texttt{gemini-3-flash-preview}; la elección se alinea con la disponibilidad de Gemini 3 Flash en el ecosistema del gateway y con los objetivos del proyecto: \textbf{baja fricción operativa}, \textbf{alto contexto} para continuidad y \textbf{costos controlables} a escala.

\section{Conclusión}
Se seleccionó \textbf{Gemini 3 Flash} por ofrecer una \textbf{ventana de contexto sustancialmente mayor} que las alternativas comparadas y un perfil de costo adecuado para una herramienta de estudio con conversaciones largas y frecuentes. En términos de arquitectura, esta decisión reduce la complejidad necesaria para gestionar historial y habilita análisis extensos con menor riesgo de pérdida de contexto, manteniendo un costo por token compatible con un esquema de adopción masiva.
